\documentclass[12pt,a4paper]{article}

\usepackage{cmap}
\usepackage[T2A]{fontenc}
\usepackage[utf8]{inputenc}
\usepackage[russian]{babel}
\usepackage{amsmath,amssymb}
\usepackage[a4paper,margin=2cm]{geometry}

\setlength{\parindent}{0pt}
\setlength{\parskip}{0.45em}
\sloppy

\newcommand{\R}{\mathbb{R}}
\newcommand{\CC}{\mathbb{C}}
\newcommand{\dd}{\,d}

\begin{document}

\begin{center}
{\Large\bfseries Рецензия на ответы к билетам}
\end{center}

\textbf{Общий вывод.}
Ответы в целом сильные: структура единообразна, определения в основном
корректны, ключевые теоремы названы, большинство доказательств доведено до
рабочих оценок. Существенных грубых математических ошибок, из-за которых
ответ надо было бы признать неверным целиком, я не обнаружил.

Главная зона роста -- не добавлять объем ради объема, а точнее расставить
гипотезы и зависимости. В устном ответе профессор будет спрашивать не только
формулу, но и почему именно эти условия нужны: где используется компактность,
где нужна равномерность, где достаточно ограниченности, а где требуется
непрерывность. Ниже -- замечания по каждому билету.

\section*{Билет 1}

\textbf{Оценка.} Ответ полный и логически устойчивый. Координатный критерий
сходимости, классификация точек множества, открытость/замкнутость и критерий
компактности изложены корректно.

\textbf{Что улучшить.}
Полезно явно сказать, что используемое определение компакта является
последовательностным, а в \(\R^n\) оно эквивалентно обычной
гейне-борелевской компактности. Это снимет возможный вопрос о том, почему
компактность доказывается через подпоследовательности.

В блоке про открытые и замкнутые множества стоит добавить одну строку с
дуальными свойствами: произвольное объединение открытых множеств открыто,
конечное пересечение открытых открыто; через дополнение получаются
соответствующие свойства замкнутых. Сейчас свойства названы, но для устного
ответа полезно явно связать их с переходом к дополнениям.

Для наглядности можно завершить определения схемой:
\[
  \overline X=X\cup X',\qquad
  \partial X=\overline X\setminus\operatorname{int}X
  =\overline X\cap\overline{\R^n\setminus X}.
\]
Это быстро показывает связь между точками прикосновения, предельными точками,
внутренностью и границей.

\section*{Билет 2}

\textbf{Оценка.} Материал раскрыт хорошо: предел по множеству, критерий
Гейне, непрерывность, свойства непрерывных функций на компакте и теорема о
промежуточных значениях даны в правильной логике.

\textbf{Что улучшить.}
В определении предела по направлению лучше явно оговорить, что здесь взят
предел по лучу \(x=x_0+t\ell\), \(t\to+0\). Если преподаватель понимает
``по направлению'' как предел вдоль всей прямой, надо сказать, что
противоположный луч соответствует направлению \(-\ell\).

В примере с одинаковыми пределами по прямым хорошо добавить одну фразу:
``проверка только прямых траекторий не исчерпывает всех способов подхода к
точке''. Это делает логический смысл примера очевидным.

В теореме о промежуточных значениях стоит подчеркнуть, что используется
линейная связность области: мы сводим задачу к одномерной функции
\(f(r(t))\) на отрезке. Без этой фразы утверждение выглядит как простое
обобщение одномерной теоремы, хотя ключевой шаг -- именно наличие кривой.

\section*{Билет 3}

\textbf{Оценка.} Определение дифференцируемости, дифференциал, необходимые и
достаточные условия, правило цепочки и инвариантность формы первого
дифференциала изложены верно. Доказательство достаточного условия через
разложение приращения по координатам хорошее.

\textbf{Что улучшить.}
Нужно явно предупредить, что существование всех частных производных в точке
само по себе не гарантирует дифференцируемость. В финальной устной версии это
сказано, но для убедительности стоит привести короткий стандартный пример или
хотя бы фразу: ``поэтому в достаточном условии важна непрерывность частных
производных в точке''.

В доказательстве производной по вектору остаток записан как \(o(t)\). Строже
писать \(o(|t|)\), а затем переходить к \(t\to+0\). Это не меняет сути, но
убирает мелкую неаккуратность.

Инвариантность формы дифференциала лучше формулировать в скалярном виде
\[
  dz=\sum_j \frac{\partial G}{\partial y_j}(y^0)\,dy_j,
\]
а затем уже упоминать матричную запись. Так студенту проще объяснить, что
``форма'' сохраняется, а смысл \(dy_j\) меняется.

\section*{Билет 4}

\textbf{Оценка.} Основной материал присутствует: высшие частные производные,
равенство смешанных производных, высшие дифференциалы и обе формы формулы
Тейлора. Доказательство формулы Тейлора через
\(\varphi(t)=f(x^0+t\Delta x)\) выбрано удачно.

\textbf{Что улучшить.}
В теореме о независимости смешанных производных лучше использовать более
чистую гипотезу \(f\in C^k\) в окрестности точки. Текущая формулировка через
существование и непрерывность производных \(k\)-го порядка в точке верна в
духе результата, но при устном доказательстве соседних перестановок индексов
легче защищать именно условие \(C^k\).

Надо четче различать объект \(d^k f(x^0)\) и его значение на приращении:
это \(k\)-линейная форма или однородный многочлен по \(dx\), а в формуле
Тейлора затем подставляют \(dx_i=\Delta x_i\). Сейчас это сказано, но можно
выделить отдельной фразой перед формулой Тейлора.

Для формы Пеано стоит добавить, что непрерывность производных до порядка
\(m\) обеспечивает не только существование \(d^m f\), но и малость разности
коэффициентов в \(d^m f(x^0+\theta\Delta x)-d^m f(x^0)\). Это ключевой
момент доказательства.

\section*{Билет 5}

\textbf{Оценка.} Ответ содержательный и в целом корректный. Определены
клеточные множества, нижняя и верхняя меры Жордана, критерий через клеточные
приближения, критерий через границу и конечная аддитивность.

\textbf{Что улучшить.}
В доказательстве критерия по границе есть места, где надо аккуратнее
обосновывать клеточность множеств вида \(B_\varepsilon\setminus
\operatorname{int}A_\varepsilon\) и \(\Pi_0\setminus
\operatorname{int}C_\varepsilon\). Для экзамена достаточно сказать, что после
общего измельчения разбиений эти разности представляются конечными
объединениями клеток.

При доказательстве операций над измеримыми множествами стоит явно назвать
факт: конечное объединение множеств меры нуль имеет меру нуль, и любое
подмножество множества меры нуль тоже имеет меру нуль. Именно это позволяет
перейти от включений для границ к измеримости.

В аддитивности лучше подчеркнуть, почему условие ``нет общих внутренних
точек'' допустимо: пересечения могут лежать на границах, а границы измеримых
по Жордану множеств имеют нулевую меру.

\section*{Билет 6}

\textbf{Оценка.} Очень полный ответ. Суммы Дарбу, критерии интегрируемости,
достаточные условия, свойства интеграла, интеграл с переменным верхним
пределом, формула Ньютона--Лейбница, замена переменной и интегрирование по
частям разобраны последовательно.

\textbf{Что улучшить.}
Теорема о среднем сформулирована со строгим условием \(g(x)\ne0\). Это
корректно, но сильнее обычной версии, где достаточно, чтобы \(g\) не меняла
знак. На экзамене стоит сказать: ``я формулирую чуть более сильную, но
достаточную для доказательства через теорему Коши версию''.

В доказательстве интегрируемости функции с конечным числом разрывов нужно
аккуратнее сказать, что вне малых окрестностей разрывов получается конечное
объединение компактных отрезков, на каждом из которых функция непрерывна, а
значит равномерно непрерывна; затем берется минимум соответствующих
\(\delta\).

Для определения интеграла Римана через суммы Дарбу стоит заранее подчеркнуть,
что речь идет об ограниченных функциях. Сейчас ограниченность названа как
необходимое условие, но в критериях она уже фактически используется.

\section*{Билет 8}

\textbf{Оценка.} Формулы геометрических приложений даны правильно, а площадь
поверхности вращения оставлена без доказательства, как требует билет.

\textbf{Что улучшить.}
В доказательстве объема тела вращения надо явно отметить, что из
интегрируемости \(f\) следует интегрируемость \(f^2\), поскольку
интегрируемые функции ограничены, а произведение интегрируемых функций
интегрируемо. Это маленькая, но важная логическая ступень.

Доказательство длины кривой лучше сделать менее ссылочным. Вместо фразы про
``достаточное условие спрямляемости'' можно вывести оценку
\[
  |r(t_i)-r(t_{i-1})|
  \le \int_{t_{i-1}}^{t_i}|r'(t)|\dd t
\]
для каждой хорды и затем просуммировать по разбиению. Так верхняя оценка
длины становится самостоятельной.

Для площади поверхности вращения полезно произнести происхождение множителя
\(2\pi f(x)\): это длина окружности радиуса \(f(x)\), умноженная на элемент
длины дуги \(\sqrt{1+(f')^2}\dd x\). Доказательство не нужно, но такая
интерпретация помогает запомнить формулу.

\section*{Билет 9}

\textbf{Оценка.} Ответ корректный: определение несобственного интеграла,
критерий Коши, сравнение, абсолютная сходимость, признаки Дирихле и Абеля
изложены в правильной последовательности.

\textbf{Что улучшить.}
В начале полезно явно разделить два типа несобственности: бесконечный
промежуток и неограниченность функции у конечной особой точки. В ответе оба
случая фактически покрыты, но для устного экзамена это лучше назвать.

В признаке Дирихле стоит подчеркнуть, что из \(g'\le0\) и \(g\to0\) следует
\(g\ge0\), поэтому
\[
  \int |g'|=g(a)-g(\beta)
\]
и появляется конечная оценка полной вариации. Сейчас это есть, но именно этот
шаг является центральным.

Для признака Абеля хорошо явно сказать, что монотонная ограниченная функция
\(g\) имеет конечный предел, после чего \(g-g_0\) переводит задачу к
Дирихле. Это делает доказательство коротким и прозрачным.

\section*{Билет 10}

\textbf{Оценка.} Ответ сильный и полный. Даны основные признаки сходимости,
преобразование Абеля, перестановка абсолютно сходящегося ряда и произведение
абсолютно сходящихся рядов.

\textbf{Что улучшить.}
В признаках Даламбера и Коши надо явно держать в голове, что они применяются
к рядам с положительными членами или к рядам из модулей. Для знакопеременных
рядов сначала переходят к \(\sum |a_k|\), если речь об абсолютной сходимости.

В признаке Лейбница лучше использовать запись \((-1)^{k-1}b_k\), если первый
член предполагается положительным. Математически текущая запись
\((-1)^k b_k\) не ошибочна, но она меняет знак суммы и может сбить на устном
ответе.

В доказательстве произведения абсолютно сходящихся рядов полезно отдельно
назвать, что речь идет не только о произведении Коши, а о любой нумерации
всех попарных произведений. Тогда становится ясно, почему перед вычислением
суммы разрешено перейти к нумерации ``квадратами''.

\section*{Билет 11}

\textbf{Оценка.} Определения равномерной сходимости, критерий Коши,
непрерывность предела, интегрирование, дифференцирование и признаки
Вейерштрасса, Дирихле и Абеля разобраны хорошо.

\textbf{Что улучшить.}
В теореме о непрерывности равномерного предела нужно сказать, что если
\(X\) не является открытым множеством, то непрерывность понимается
относительно \(X\). Это убирает возможную претензию к \(\delta\)-окрестностям
на произвольном множестве.

В теореме о дифференцировании важно подчеркнуть необходимость условия
сходимости \(f_n(x_0)\) хотя бы в одной точке. Одна равномерная сходимость
производных не фиксирует аддитивную константу.

В признаке Дирихле для функциональных рядов надо аккуратно формулировать
монотонность \(b_k(x)\): она требуется по \(k\) при каждом фиксированном
\(x\), а стремление к нулю -- равномерное по \(x\). Именно сочетание этих
двух условий дает равномерную оценку хвоста.

\section*{Билет 12}

\textbf{Оценка.} Ответ корректный: радиус сходимости, формула
Коши--Адамара, характер сходимости внутри и вне круга, равномерная
сходимость внутри круга, первая теорема Абеля и неизменность радиуса при
формальном дифференцировании/интегрировании изложены верно.

\textbf{Что улучшить.}
В формуле Коши--Адамара стоит проговорить особые случаи:
если \(\limsup\sqrt[k]{|c_k|}=0\), то \(R=+\infty\); если limsup бесконечен,
то \(R=0\). Это часто спрашивают отдельно.

В доказательстве неизменности радиуса для производного ряда лучше сразу
сравнивать коэффициенты ряда \(\sum_{k=1}^\infty kc_k z^{k-1}\) после
переиндексации. Текущий путь через \(\sum kc_k z^k\) верен, но на слух
требует лишнего пояснения о множителе \(z\).

На границе круга стоит подчеркнуть, что радиус сходимости ничего не говорит
о конкретных граничных точках: возможна сходимость во всех, ни в одной или
только в части точек.

\section*{Билет 13}

\textbf{Оценка.} Ответ широкий и в основном правильный. Хорошо раскрыты
почленное интегрирование и дифференцирование, единственность коэффициентов,
ряд Тейлора, интегральный остаток, пример гладкой неаналитической функции и
основные разложения.

\textbf{Что улучшить.}
Достаточное условие регулярности через одну общую константу
\(|f^{(n)}(x)|\le M\) очень сильное. Оно верное, но стоит подчеркнуть, что
это лишь удобное достаточное условие, а не критерий разложимости. Более общая
идея: надо доказать \(r_n(x)\to0\) в формуле Тейлора.

Для разложения \(\ln(1+x)\) при \(x=1\) используется ссылка на теорему
Абеля. Ее надо либо кратко сформулировать, либо в устном ответе оставить
сначала интервал \((-1,1)\), а затем отдельно сказать, что точка \(x=1\)
добавляется по предельному переходу.

В биномиальном разложении полезно отдельно оговорить, что поведение в
точках \(x=\pm1\) зависит от \(\alpha\); поэтому безопасная основная область
для произвольного \(\alpha\) -- \((-1,1)\).

\section*{Билет 14}

\textbf{Оценка.} Ответ правильный для безусловных локальных экстремумов.
Необходимое условие через теорему Ферма и достаточные условия через знак
второго дифференциала доказаны корректно.

\textbf{Что улучшить.}
Так как в определениях упомянут условный экстремум, надо явно сказать, что
дальнейшие необходимые и достаточные условия относятся к безусловному
экстремуму во внутренней точке. Иначе может возникнуть вопрос, почему не
появляются ограничения и множители Лагранжа.

Для практической проверки знака квадратичной формы стоит добавить критерий
Сильвестра для матрицы Гессе: положительная определенность эквивалентна
положительности всех ведущих главных миноров, отрицательная -- чередованию
их знаков. Это не обязательно для доказательства, но очень полезно в ответе.

В вырожденном случае лучше говорить не только ``полуопределенная форма'', а
точнее: если второй дифференциал не является ни положительно определенным, ни
отрицательно определенным, ни знаконеопределенным, тест второго порядка может
не решить вопрос. Примеры \(x^4\) и \(x^3\) хорошо оставлять именно здесь.

\section*{Итоговая рекомендация}

Перед экзаменом я бы не переписывал ответы заново. Рациональнее пройтись по
ним как по конспекту и на полях добавить три типа пометок:

1. Где используется ключевая гипотеза: компактность, равномерность,
монотонность, ограниченность, непрерывность производных.

2. Где формулировка дана в более сильной версии, чем минимально нужна:
теорема о среднем с \(g\ne0\), интегрирование/дифференцирование рядов с
непрерывными членами, достаточное условие регулярности.

3. Где полезен один короткий пример: несуществование предела при одинаковых
пределах по прямым, частные производные без дифференцируемости, стационарная
точка без экстремума, гладкая неаналитическая функция.

Такая доработка повысит ясность ответа больше, чем добавление новых длинных
доказательств.

\end{document}
